% Generated draft: requires mathematical and source review.
\documentclass[11pt]{article}
\usepackage[margin=25mm]{geometry}
\usepackage{fontspec}
\setmainfont{Latin Modern Roman}
\usepackage{amsmath,amssymb,mathtools,unicode-math}
\setmathfont{Latin Modern Math}
\usepackage{enumitem,needspace,xurl,hyperref}
\setlength{\parindent}{0pt}
\setlength{\parskip}{5pt}
\setlength{\emergencystretch}{3em}
\newcommand{\N}{\mathbb{N}}
\newcommand{\Z}{\mathbb{Z}}
\newcommand{\Q}{\mathbb{Q}}
\newcommand{\R}{\mathbb{R}}
\newcommand{\C}{\mathbb{C}}
\newcommand{\F}{\mathbb{F}}
\newcommand{\A}{\mathbb{A}}
\newcommand{\PP}{\mathbb{P}}
\newcommand{\E}{\mathbb{E}}
\newcommand{\eps}{\varepsilon}
\newcommand{\dd}{\,\mathrm{d}}
\begin{document}
\section*{UnsolvedMath ID: 30005453}
\textbf{Source problem code:} catalogue entry
\subsection{Definitions and mathematical statement}
Let \(G=(V,E)\) be an infinite graph of bounded degree, so that \(\sup_{v\in V}\deg_G(v)<\infty\). Write \(p(v)\) for the firing rate at vertex \(v\), and assume that these rates are bounded above. Let \(\alpha<1\) be the reinforcement parameter, and write \(X_e(t)\) for the source's normalized reinforcement variable associated with edge \(e\in E\) at time \(t\).

The supplied record does not define the reinforcement process, its normalization, its initial conditions, or the equations and admissible state space defining an equilibrium. Consequently, the following is a formulation of the intended conjecture conditional on those missing definitions, rather than a complete standalone specification of a stochastic model. No additional lower bound on \(\alpha\), positivity or uniform lower bound for \(p(v)\), or graph assumptions beyond those stated above are imposed here.

\textbf{Subcritical conjecture, as supplied.} For every graph, firing-rate assignment, and parameter satisfying the preceding hypotheses, does the intended reinforcement model have exactly one admissible equilibrium \(x=(x_e)_{e\in E}\), with
\[
\text{for every }e\in E,\qquad
\mathbb P\!\left(\lim_{t\to\infty}X_e(t)=x_e\right)=1?
\]
The probability law and the quantification over initial conditions remain unspecified in the supplied record. If \(E\) is countable, the displayed assertion also gives convergence for all edges on a single event of probability one. Countability is not explicitly stated in the record.

\textbf{Supplied subcritical context.} The excerpt asserts that, for finite graphs, concavity of a function \(L\) gives convergence to a unique equilibrium maximizing \(L\) on a nonnegative simplex. Neither \(L\) nor a reliably typeset normalization constraint is supplied. It attributes partial results to Y. Couzini\'e and C. Hirsch, reference [2, Theorem 2.2], under the assumption \(p(v)\equiv1\): convergence for \(\alpha<\tfrac12\), and convergence for \(\alpha<1\) when \(G=\mathbb Z\). The supplied dated assessment describes the first result as applying to bounded-degree graphs. These are reported source claims, not independently verified results here.

\textbf{Separate critical context retained from the excerpt.} The source also discusses \(\alpha=1\), which is outside the subcritical conjecture. It describes \(L\) as non-strictly concave and reports nonuniqueness of equilibrium on even cycles, citing reference [8]. For uniform firing rates, it gives the equilibria on a cycle of length \(2k\) in the form
\[
\{(a,b,a,b,\ldots,a,b):a+b=2\},
\]
without explicitly specifying the admissible domains of \(a,b\). It reports uniqueness on odd cycles and claims that, as \(k\to\infty\), the law of \((a,b)\) tends to the Dirac measure at \((1,1)\); the law selecting these equilibria is not defined in the record. This motivates a separate critical conjecture for \(G=\mathbb Z\) and \(p(v)\equiv1\), rendered in the excerpt as almost surely \(x_v(t)\to1\) for every \(v\). The relation of that notation to the edge variables \(X_e(t)\) is unspecified. This critical conjecture is retained as source context and is not incorporated into the subcritical target.
\subsection{Short English statement}
The supplied subcritical conjecture asks whether every infinite bounded-degree graph with bounded firing rates has a unique equilibrium and almost-sure convergence of every normalized edge variable to that equilibrium when \(\alpha<1\). The process, equilibrium equations, and initial-condition assumptions are missing, so a complete standalone mathematical statement requires further source information. The supplied partial results and separate critical-case discussion are retained as attributed context.
\subsection{Sources}
\begin{itemize}
\item Supplied local registry record, permanent UnsolvedMath ID \(30005453\), exact source problem code \texttt{catalogue entry}, titled ``Unique Equilibria for Subcritical Reinforcement on Infinite Graphs.'' No individual proposer is identified; the recorded proposed year is 2023.
\item \emph{MATRIX-MFO Tandem Workshop: Stochastic Reinforcement Processes and Graphs} (2023), Oberwolfach Reports, pp. 639--679, DOI 10.4171/owr/2023/12. Supplied source and literature link: \url{https://doi.org/10.4171/owr/2023/12}.
\item Couzini\'e--Hirsch, ``weakly reinforced P\'olya urns,'' supplied literature link: \url{https://arxiv.org/abs/2010.03347}. The excerpt attributes the stated subcritical partial results to Y. Couzini\'e and C. Hirsch through reference [2, Theorem 2.2]; the complete report bibliography is not supplied.
\item The critical-cycle discussion cites reference [8], whose bibliographic identity is not supplied.
\item Literature assessment checked at \texttt{2026-08-22}: the supplied curation labels the problem ``partially solved'' and reports convergence below \(\alpha=\tfrac12\) on bounded-degree graphs and below \(\alpha=1\) on \(\mathbb Z\), while describing unique-equilibrium convergence on arbitrary infinite graphs as remaining open. This is supplied dated curation, not an independently verified literature assessment; the excerpt states \(p(v)\equiv1\) for the attributed partial results.
\item The registry supplies the statement-status label \texttt{corrected\_verified} and says its clean statement was checked against the official report. That verification claim is part of the supplied record and has not been independently checked here.
\end{itemize}
\end{document}
